\section{Conclusão}\label{sec:conclusao}

Este trabalho perguntou em que medida as transferências intergovernamentais condicionam o financiamento e o gasto de Manaus e Parintins entre 2016 e 2025 e se esse padrão é compatível com o efeito \textit{flypaper} e com a ilusão fiscal, à luz da organização fiscal subnacional da Espanha. A resposta é que as transferências condicionam fortemente ambos os municípios, em intensidades muito diferentes. Parintins financiou, em média, \depParMedia\% de sua receita com transferências, e sua receita tributária correspondeu a apenas \autParMedia\% da receita corrente líquida; Manaus, sustentada pela base econômica da Zona Franca, dependeu de transferências em \depManMedia\% da receita, com autonomia tributária crescente. Nas regressões com os dados anuais dos dois municípios, a despesa acompanha as transferências --- R\$ \bTrUm{} por real transferido --- e responde pouco à renda local --- R\$ \bPibUm{} por real de PIB per capita ---, diferença estatisticamente significativa na especificação principal e compatível com o efeito \textit{flypaper}. Esse padrão, porém, vem sobretudo de Parintins, onde a despesa cresce quase real a real com as transferências (\bTrDez{} por real), e perde significância quando se controla a tendência comum ou se usam primeiras diferenças; com tão poucas observações, ele deve ser lido como evidência sugestiva, e não conclusiva. Não se encontrou, contudo, associação entre transferências e arrecadação tributária, o que não sustenta a versão da ilusão fiscal baseada na substituição de esforço tributário.

Em relação às hipóteses, H1 (dependência maior em Parintins do que em Manaus e nos municípios do que nos subsetores espanhóis) foi confirmada para Parintins e, quanto à autonomia tributária, também para Manaus; H2 (propensão a gastar as transferências maior que a propensão a gastar a renda) foi confirmada parcialmente, de forma nítida em Parintins e na especificação em níveis, mas não nas especificações com tendência e em primeiras diferenças; H3 (redução do esforço tributário) não foi confirmada; e H4 (maior disciplina dos níveis subnacionais espanhóis do que da administração central) foi confirmada pelos dados do Eurostat, com Comunidades Autônomas próximas do equilíbrio e entidades locais superavitárias em nove de dez anos.

As contribuições do trabalho são três. A primeira é uma base de dados reprodutível, que preserva o trabalho anterior, documenta cada fonte e cada transformação e corrige problemas que afetavam diretamente as conclusões, como a mistura de bases populacionais, a inclusão de operações intraorçamentárias e o uso de um conceito inadequado de resultado para a Espanha. A segunda é empírica: o estudo traz para a literatura brasileira sobre efeito \textit{flypaper} dois municípios da Região Norte, pouco estudados, e mostra que, no município mais dependente, a associação entre transferências e gasto opera pela expansão da despesa, e não pela redução de tributos. A terceira é comparativa: o uso de dados por subsetor das contas nacionais espanholas permite comparar níveis de governo semelhantes, em vez de comparar municípios com o governo geral de um país.

As implicações para as finanças públicas são diretas. Em municípios com base tributária ínfima, como Parintins, a qualidade do gasto depende quase integralmente do desenho das transferências: regras de vinculação, critérios de rateio e mecanismos de prestação de contas são as principais alavancas de política. A transição para o IBS deve ser acompanhada de perto, porque redistribuirá justamente as receitas (ISS e cota-parte do ICMS) que diferenciam Manaus de Parintins. E a experiência espanhola sugere que ampliar a responsabilidade dos governos subnacionais pela arrecadação, combinada com regras fiscais críveis, é compatível com descentralização ampla.

As conclusões respeitam o alcance dos dados e do método. As estimativas medem associações condicionais com apenas 15 observações, e não efeitos causais; o PIB municipal é uma \textit{proxy} imperfeita da renda local e termina em 2023; a população de 2023 é interpolada; Parintins não tem DCA de 2016 no Siconfi; e a comparação com a Espanha é qualitativa, dadas as diferenças de classificação contábil e de competências. Pesquisas futuras podem ampliar a amostra para outros municípios e explorar as descontinuidades das faixas do FPM, com estimadores robustos à heterogeneidade dos efeitos \cite{callaway2021,dechaisemartin2020}, analisar separadamente transferências vinculadas e livres, incorporar indicadores de resultado das políticas (como o IDEB e indicadores de saúde) e acompanhar os efeitos da Reforma Tributária sobre a autonomia dos municípios amazônicos.
